\documentclass[10pt,a4paper]{article}
\usepackage{amsmath,amssymb,amsthm,graphicx}
\usepackage{hyperref}
\usepackage{amsmath}
\usepackage{enumerate}
\usepackage{amsfonts}
\usepackage{amsthm}
\newtheorem{theorem}{Theorem}
\newtheorem{proposition}[theorem]{Proposition}
\newtheorem{lemma}[theorem]{Lemma}
\newtheorem{corollary}[theorem]{Corollary}
\newtheorem{rem
ark}[theorem]{Remark}
\newtheorem{definition}[theorem]{Definition}
\newtheorem{example}[theorem]{Example}
\newcommand{\out}{\textup{out}} 
\newcommand{\spa}{\textup{span}} 
\newcommand{\Stab}{\textup{Stab}} 
\newcommand{\peb}{\textup{peb}} 

\newcommand{\AcceptEdge}[1]{\textbf{Accept-Edge}(#1)}
\newcommand{\MovePebble}[1]{\textbf{Move-Pebble}(#1)}
\title{Sparsity is unchanged when all edges live on fewer than n vertex classes}
\author{Signe Lundqvist, Tovohery Randrianarisoa, Klara Stokes, Joannes Vermant}
\date{arxiv:2306.05050v3; Creative Commons Attribution 4.0 International}
\begin{document}
\maketitle
\noindent\textit{Editorial note.} The counted statement is the author's unlabelled lemma environment at source0.tex lines 289-293; its proof environment opens at line 294, immediately after \textbackslash{}end\{lemma\}, so the pairing is by adjacency and the retained proof range 294-308 is the whole environment.\par\medskip
\noindent\textit{Editorial note.} The numbered system of inequalities at source lines 241-245 is retained as context because the statement assumes it and the proof applies it at source lines 300 and 307; it is the only outside label of the retained proof.\par\medskip
\noindent\textit{Editorial note.} The retained proof carries no citation, so this card needs no bibliography. The source does ship both a bibliography.bib and a prepared retained bibliography, which the cards of this source that do cite will use.\par\medskip
\noindent\textit{Editorial note.} Preamble: the author's own package lines, theorem-environment family and macro block (source lines 4-39); the drawing and algorithm packages of the source are not retained because the retained material does not use them. In this source the theorem declarations stand after the beginning of the document, which does not affect their meaning.\par\medskip
\noindent\textit{Editorial note.} Metadata provenance: the licence evidence file records source\_version 2306.05050v3 and status supported, and the card's source\_version is that value; the author list is taken from the authors\_index key of the source record file of the pinned copy, transcribed from its LaTeX escapes; the title is the source record's title key.\par\medskip
\noindent\textit{Editorial note.} The source carries no label-number dump and no counter values are imposed, so the printed numbering is the wrapper's default; identity is the pinned source position, the author having given no label.\par\medskip
\noindent\textit{Editorial scope.} One original result of the pinned author manuscript, counted once: sparsity is preserved when the edges are supported on a proper initial set of vertex classes (unlabelled lemma at source0.tex lines 289-293).\par\medskip
\noindent\textit{Allowed standard prerequisites (editorial).}\begin{itemize}
\item Sparsity parameters. A pair of integer sequences K and L, together with an ordered partition of the vertex set into classes, so that every subset of vertices has its edge count bounded by the weighted size of the classes it meets minus the corresponding parameter.. It is the notion of sparsity under consideration.
\item Support of a graph. The edges of the graph are supported on an initial run of the vertex classes when no edge meets a later class.. It is the hypothesis of the statement.
\item The numbered parameter inequalities. Three inequalities relating consecutive parameter entries, the first ensuring that a single edge is sparse and the others ensuring that sparsity survives adding isolated vertices.. They are the hypotheses invoked in the proof.
\end{itemize}\medskip
\begin{align}
    k_i + k_{i+1} &> l_{i, i+1} \label{conditionposet1 -k k l},\\
    k_{i+1} + l_{i-1, i} &> l_{i, i+1}\label{conditionposet2 - k l l},\\
     l_{1,2} \leq  \dots &\leq l_{n-1,n}. \label{conditionposet3-increasing}
\end{align}

\begin{lemma}
    Let $K=(k_1, \dots, k_n)\in  \mathbb{Z}_{>0}^{n}$ and $L= (l_{1,2}, \dots, l_{n-1,n})\in  \mathbb{Z}_{\geq 0}^{n-1}$ be such that $K$ and $L$ satisfy inequalities \eqref{conditionposet1 -k k l}, \eqref{conditionposet2 - k l l}, and \eqref{conditionposet3-increasing}. Let $G=(V,E)$ be a (multi)graph such that $V$ has a partition into subsets $V_1, V_2, \dots, V_n$ such that for all $e\in E$ there is some $i$ such that $e \in E(V_i, V_{i+1})$.
    
         Suppose that all edges are supported on $V_1\cup \dots \cup V_{t}$ with $t<n$. Then $G$ is $(K, L)$-sparse if and only if $G'= (V_1\cup \dots \cup V_{t}, E)$ is $(K',L')$-sparse, where $K'= (k_1, \dots, k_t)$ and $L'= (l_{1,2}, \dots, l_{t-1, t})$.
\end{lemma}

\begin{proof}
`Only if' is clear. Thus, suppose that $G'$ is a $(K',L')$-sparse graph. It suffices to show that sparsity is preserved when adding a set $W$ of isolated vertices of $V_{t+1}, \dots, V_n$ to a set $V' \subseteq V_{1}\cup \dots \cup V_t$. Assume first $W\cup V' \subseteq V_1\cup \dots \cup V_{t+1}$, with $V' \cap V_{t-1} \neq \emptyset$, $ V' \cap V_{t}\neq \emptyset$, and $W \cap V_{t+1} \neq \emptyset$. Then inequality \eqref{conditionposet2 - k l l} and the sparsity of $G'$ imply that
\begin{align*}
   \vert E(W\cup V') \vert&= \vert E(V') \vert \leq \sum_{i=1}^{t} k_i \vert V'\cap V_i\vert - l_{t-1,t}\\
&< \sum_{i=1}^{t} k_i \vert V' \cap V_i\vert + k_{t+1} - l_{t, t+1} \leq \sum_{i=1}^{t+1} k_i \vert (W\cup V') \cap V_i\vert - l_{t, t+1}, 
\end{align*}
and hence, since inequality \eqref{conditionposet3-increasing} ensures that $l(V' \cup W)=l_{t, t+1}$ in the last expression, this subset meets the sparsity condition.
If $V' \cup W$ is a set with $W \cap V_{t+j} \neq \emptyset$, $W \cap V_{t+j+1}\neq \emptyset$ for some $j > 0$, then by taking the largest such $j$, we see that 
\begin{align*}
\vert E(W\cup V') \vert&= \vert E(V') \vert \leq \sum_{i=1}^{t} k_i \vert V '\cap V_i\vert - l(V')\\
&<  \sum_{i=1}^{t} k_i \vert V '\cap V_i\vert - l(V') + (k_{t+j} + k_{t+j+1} - l_{t+j, t+j+1}) \\
&\leq \sum_{i=1}^{n} k_i \vert (V'\cup W) \cap V_i\vert - l(V'\cup W)
\end{align*}
where we used inequality \eqref{conditionposet1 -k k l} to go from the first line to the second. For the case where $(V'\cup W) \cap V_{t} \neq \emptyset$, $W \cap V_{t+1}\neq \emptyset$, and $V_{t-1}\cap V'= \emptyset$, one similarly uses  \eqref{conditionposet1 -k k l}, and in any other case, it is trivial to show that $V'\cup W$ satisfies the sparsity condition. 
\end{proof}

\end{document}
